\documentclass[11pt]{article}

\usepackage[margin=1in]{geometry}
\usepackage[T1]{fontenc}
\usepackage{amsmath,amssymb,bm}
\usepackage{booktabs}
\usepackage{array}
\usepackage{xcolor}
\usepackage{hyperref}

\hypersetup{
  colorlinks=true,
  linkcolor=blue!55!black,
  urlcolor=blue!55!black,
  citecolor=blue!55!black
}

\newcommand{\code}[1]{\texttt{\detokenize{#1}}}
\newcommand{\R}{\mathbb{R}}
\newcommand{\C}{\mathbb{C}}
\newcommand{\ii}{\mathrm{i}}
\newcommand{\dd}{\mathrm{d}}
\newcommand{\norm}[1]{\left\lVert #1 \right\rVert}
\newcommand{\angles}{N_{\Omega}}
\newcommand{\om}{\Omega}
\newcommand{\Ah}{A_h}
\newcommand{\Bh}{B_h}
\newcommand{\Geta}{G_\eta}
\newcommand{\QK}{Q_K}

\title{Base and Learned NPBS Preconditioners for the Periodic RTE Prototype}
\author{Project note assembled from repository audits and ablation discussions}
\date{May 26, 2026}

\begin{document}
\maketitle

\begin{abstract}
This note records the mathematical model, discretization, classical
preconditioners, learned variants, training losses, diagnostics, and current
experimental conclusions for the radiative-transfer-equation (RTE) part of the
Neural Preconditioned Born Series (NPBS) repository.  It is intended to be
self-contained enough that a new researcher can continue the small-scale laptop
experiments without reconstructing the notation from code and chat history.

The current RTE prototype is a periodic two-dimensional discrete-ordinates
model with pseudo-spectral spatial differentiation.  The strongest learned
result so far is a full-angular \code{UNet-defect2} model trained by physical
one-step residual loss: it can nearly match a deeper classical full-TSA
correction on the hard mixed case.  The structured moment-space models
(\code{CoeffWeightedTSA}, \code{CoeffMixedTSA}, \code{LearnedRP-TSA}, and
\code{ResidualSubspaceTSA}) show that mode coupling and learned
restriction/prolongation can reduce one-step residuals, but the same learned
maps can damage FGMRES convergence.  The most important next direction is a
stabilized learned angular subspace architecture: learn restriction and
prolongation from a larger candidate angular space, preserve the scalar mode,
control the operator norm, and checkpoint by Krylov behavior rather than only
by one-step residual loss.
\end{abstract}

\section{Repository map}

The RTE code path used in the NPBS experiments is concentrated in the following
files.

\begin{center}
\begin{tabular}{>{\raggedright\arraybackslash}p{0.34\linewidth}
                >{\raggedright\arraybackslash}p{0.58\linewidth}}
\toprule
File & Role \\
\midrule
\path{rte/models.py}
& Defines \code{PeriodicRTE2D}: the periodic discrete-ordinates operator,
FFT streaming, reference inverse, DSA/TSA synthetic corrections, Richardson,
and FGMRES-facing preconditioner application. \\
\path{rte/run_hard_case_experiments.py}
& Defines the benchmark coefficient fields and sources for the four hard
cases: \code{diffusive}, \code{heterogeneous}, \code{anisotropic}, and
\code{mixed}. \\
\path{rte/npbs_unet_ablation.py}
& Main ablation runner for U-Net and moment-space learned preconditioners;
contains \code{UNet-residual}, \code{UNet-defect2}, \code{MomentUNet},
\code{CoeffWeightedTSA}, \code{CoeffMixedTSA}, \code{LearnedRP-TSA}, and
\code{ResidualSubspaceTSA}. \\
\path{rte/npbs_bottleneck_harness.py}
& Targeted harness for bottleneck experiments: loss modes, teacher terms,
identity anchoring, checkpointing by heldout metrics, and mode-expressivity
diagnostics. \\
\path{rte/NPBS_MODE_SUBSPACE_ARCHITECTURE_ASSESSMENT.md}
& Current narrative assessment of mode/subspace architectures and checkpoint
experiments. \\
\bottomrule
\end{tabular}
\end{center}

The older \path{rte/mgRTE_ablation} path contains experiments for the attached
ATFPS-style multigrid model.  The present note intentionally focuses on the
NPBS-native periodic pseudo-spectral prototype, because the current research
question is about a learned preconditioner based on NPBS rather than a separate
ATFPS discretization.

\section{Continuous RTE model}

The steady mono-energetic transport equation in two spatial dimensions is
written as
\begin{equation}
  \om\cdot \nabla_x \psi(x,\om)
  + \sigma_t(x)\psi(x,\om)
  - \sigma_s(x)\int_{\mathbb{S}^1} p(\om,\om')\psi(x,\om')\,\dd\om'
  = q(x,\om),
  \label{eq:continuous-rte}
\end{equation}
where \(x\in [0,1)^2\) is periodic in the current prototype,
\(\om\in\mathbb{S}^1\) is the angular direction, \(\sigma_t>0\) is total
cross section, \(\sigma_s\ge 0\) is scattering cross section, and
\(\sigma_s<\sigma_t\) is enforced in the verifier.  The scalar flux is
\begin{equation}
  \phi(x)=\int_{\mathbb{S}^1}\psi(x,\om)\,\dd\om.
\end{equation}
The hard regimes for RTE preconditioning are not only high contrast.  The
difficult cases are produced by combinations of
\begin{itemize}
  \item high albedo, where \(\sigma_s/\sigma_t \approx 1\) and the system has
  slow diffusive error modes;
  \item forward-peaked anisotropic scattering, where low scalar moments are not
  enough to represent the important angular error;
  \item mixed optically thin and optically thick regions, where streaming and
  diffusion coexist and one local acceleration mechanism is not uniformly
  correct;
  \item strong coefficient variation, where a constant reference inverse can be
  a poor global model;
  \item non-normality, where a one-step residual contraction can coexist with
  poor Krylov convergence.
\end{itemize}

\section{Discrete ordinates and periodic pseudo-spectral discretization}

The code uses uniform angular quadrature with \(\angles\) directions
\begin{equation}
  \theta_\ell = \frac{2\pi \ell}{\angles},\qquad
  \om_\ell=(\cos\theta_\ell,\sin\theta_\ell),\qquad
  w_\ell=\frac{1}{\angles},
  \qquad \ell=0,\ldots,\angles-1.
\end{equation}
The unknown is stored as
\[
  \psi_h \in \C^{B\times \angles\times n_x\times n_y},
\]
with batch dimension \(B\).  The discrete scattering term is
\begin{equation}
  (\mathcal{S}_h\psi)_\ell(x_i)
  =
  \sigma_s(x_i)\sum_{m=0}^{\angles-1}w_m p_{\ell m}\psi_m(x_i).
\end{equation}
For \(g=0\), \(p_{\ell m}=1\), giving isotropic scattering.  For \(g\ne 0\),
the code uses a circular Henyey--Greenstein-like kernel
\begin{equation}
  \tilde p_{\ell m}
  =
  \frac{1-g^2}{1-2g\cos(\theta_\ell-\theta_m)+g^2},
  \qquad
  p_{\ell m}
  =
  \frac{\tilde p_{\ell m}}
       {\sum_m w_m\tilde p_{\ell m}},
\end{equation}
so every angular row is quadrature-normalized.

On the periodic grid with spacing \(h=1/n_x=1/n_y\), spatial derivatives are
applied spectrally.  If \(\widehat{\psi}_\ell(k)\) denotes the two-dimensional
FFT coefficient, then
\begin{equation}
  \widehat{\om_\ell\cdot\nabla_h\psi_\ell}(k)
  =
  \ii(\om_{\ell,x}k_x+\om_{\ell,y}k_y)\widehat{\psi}_\ell(k),
\end{equation}
where \(k_x,k_y\) are generated by \code{torch.fft.fftfreq}.  Therefore the
discrete system is
\begin{equation}
  \Ah \psi = q,
  \qquad
  (\Ah\psi)_\ell
  =
  \om_\ell\cdot\nabla_h \psi_\ell
  +\sigma_t\psi_\ell
  -\sigma_s\sum_m w_m p_{\ell m}\psi_m .
  \label{eq:discrete-A}
\end{equation}

This is a convenient NPBS testbed, but it is not yet a production transport
discretization with inflow/vacuum boundaries, upwinding, or positivity
preservation.  Conclusions about learned preconditioner structure are useful,
but conclusions about production RTE boundary treatment require a separate
boundary-compatible discretization.

\section{Reference split and Born identity}

The reference transport operator is
\begin{equation}
  L_{\eta,\ell}
  =
  \om_\ell\cdot\nabla_h + \sigma_0,
  \qquad
  \sigma_0=\sigma_{\rm ref}+\eta.
\end{equation}
In the code,
\begin{equation}
  \sigma_{\rm ref}
  =
  \begin{cases}
    \operatorname{mean}_x\sigma_t(x), & \text{default},\\
    \text{user supplied constant}, & \text{optional},
  \end{cases}
\end{equation}
and
\begin{equation}
  \eta
  =
  \eta_{\rm floor}
  +
  \eta_{\rm scale}\,
  \max_x\left(|\sigma_t(x)-\sigma_{\rm ref}|+|\sigma_s(x)|\right).
\end{equation}
The inverse \(\Geta=L_\eta^{-1}\) is diagonal in spatial Fourier variables:
\begin{equation}
  \widehat{(\Geta r)_\ell}(k)
  =
  \frac{\widehat r_\ell(k)}
       {\ii(\om_{\ell,x}k_x+\om_{\ell,y}k_y)+\sigma_0}.
  \label{eq:Geta}
\end{equation}
Define \(V_\eta=L_\eta-\Ah\).  Then
\begin{equation}
  \Geta\Ah
  =
  I-\Geta V_\eta.
\end{equation}
The classical Born/Richardson error propagator is \(I-\Geta\Ah=\Geta V_\eta\).
The preconditioning question is how to replace \(\Geta\) by a stronger but
still cheap \(\Bh\) so that \(I-\Ah\Bh\) is small enough for stationary
iteration or has a favorable spectrum for FGMRES.

\section{Classical base preconditioners}

\subsection{Transport/Born preconditioner}

The simplest NPBS/RTE preconditioner is the shifted reference transport inverse
\begin{equation}
  B_{\rm tr}r=\Geta r.
\end{equation}
This is the direct analogue of using the shifted background Green operator in
the Helmholtz NPBS/Born method.  It is matrix-free and FFT-applied.

\subsection{DSA-style synthetic acceleration}

Diffusion synthetic acceleration (DSA) adds a low-order diffusion correction to
the residual left by the reference transport step:
\begin{align}
  z_t &= \Geta r,\\
  s &= r-\Ah z_t,\\
  \left(-D_{\rm ref}\Delta_h+\sigma_{a,\rm ref}\right)\delta\phi
    &= \Phi(s),\\
  B_{\rm DSA}r &= z_t+\mathcal{L}(\delta\phi).
\end{align}
Here \(\Phi(s)=\sum_\ell w_\ell s_\ell\) is scalar angular averaging, and
\(\mathcal{L}\) lifts a scalar correction back to the angular unknown.  In the
current periodic prototype,
\[
  D_{\rm ref}=\frac{1}{2\sigma_{\rm ref}},
  \qquad
  \sigma_{a,\rm ref}=\sigma_{\rm ref}-\sigma_{s,\rm ref},
\]
with coefficient floors for numerical safety.  The diffusion solve is also
periodic and FFT-based.  This DSA is therefore a synthetic low-order
preconditioner in the periodic NPBS testbed; it is not the full transport DSA
theory for upwind/inflow boundary discretizations.

\subsection{TSA-style synthetic acceleration}

Transport synthetic acceleration (TSA) replaces the diffusion correction by a
reference transport-scattering solve.  Let \(A_{\rm ref}\) be the constant
coefficient reference transport-scattering operator,
\begin{equation}
  (A_{\rm ref}u)_\ell
  =
  \om_\ell\cdot\nabla_h u_\ell
  +\sigma_{\rm ref}u_\ell
  -\sigma_{s,\rm ref}\sum_m w_m p_{\ell m}u_m .
\end{equation}
The full-angular TSA correction is
\begin{equation}
  B_{\rm TSA}r
  =
  \Geta r + A_{\rm ref}^{-1}(r-\Ah\Geta r).
  \label{eq:full-tsa}
\end{equation}
Because \(A_{\rm ref}\) has constant spatial coefficients and periodic
boundaries, it is block diagonal in spatial Fourier variables.  The code can
therefore apply \(A_{\rm ref}^{-1}\) by small dense angular solves per Fourier
mode.

\subsection{Moment TSA}

Moment TSA uses angular Fourier modes
\[
  \operatorname{range}(\QK)=\operatorname{span}\{0,\pm 1,\ldots,\pm K\}.
\]
The correction is
\begin{equation}
  B_{{\rm MTSA},K}r
  =
  \Geta r
  +
  \QK(\QK^*A_{\rm ref}\QK)^{-1}\QK^*(r-\Ah\Geta r).
  \label{eq:moment-tsa}
\end{equation}
This is cheaper than full TSA and is the natural base for learned angular
subspace experiments.

\paragraph{Important naming convention.}
The repository uses two different notions of ``4'':
\begin{itemize}
  \item \code{tsa4}, \code{dsa4}, or a baseline depth of \(4\) means four
  residual-correction applications of the selected base method:
  \(z\leftarrow z+\omega B_0(r-\Ah z)\) repeated four times.
  \item \code{TSA-M4}, \code{tsa_moment4}, or \code{moment_harmonics=4} means
  the angular Fourier moment space \(\{0,\pm1,\ldots,\pm4\}\).
\end{itemize}
These are not the same.  Reports should always state whether the suffix is a
depth count or an angular harmonic cutoff.

\subsection{Equal-cost comparisons}

In the ablation reports, ``equal cost'' means equal counts of selected base
preconditioner applications.  This count includes repeated calls to
\(B_0\), but it does not fully price all implementation overheads such as
\(\Ah\) applications, FFT setup, reductions, U-Net evaluation, metric
evaluation, Python overhead, or dense angular micro-solves.  Equal base-call
cost is therefore a useful algorithmic normalization, but wall-clock time and
FGMRES residual histories must still be reported.

\section{Hard cases and sources}

The hard cases are generated in \path{rte/run_hard_case_experiments.py}.  The
default grid is periodic on \([0,1)^2\) with \(h=1/n\).

\begin{center}
\begin{tabular}{>{\raggedright\arraybackslash}p{0.18\linewidth}
                >{\raggedright\arraybackslash}p{0.24\linewidth}
                >{\raggedright\arraybackslash}p{0.43\linewidth}}
\toprule
Case & Key coefficients & Purpose \\
\midrule
\code{diffusive}
& \(\sigma_t=20\), albedo \(0.99\), \(g=0\).
& Homogeneous optically thick high-albedo regime.  Tests diffusion-like slow
error modes. \\
\code{heterogeneous}
& Inclusion/channel contrast; albedo roughly \(0.82\) to \(0.97\), \(g=0\).
& Tests coefficient heterogeneity without anisotropic scattering. \\
\code{anisotropic}
& \(\sigma_t=5\), albedo \(0.96\), \(g=0.85\), beamed source.
& Tests forward-peaked scattering and angular structure. \\
\code{mixed}
& Thin left region with \(\sigma_t\approx 0.35\), thick right region with
\(\sigma_t\approx16\), albedo from \(0.20\) to \(0.985\), \(g=0.65\).
& Hardest laptop case: streaming and diffusive regimes are coupled across a
sharp interface. \\
\bottomrule
\end{tabular}
\end{center}

Training and heldout evaluations use deterministic Gaussian source families in
the current scripts.  The heldout source is not an independent coefficient
field; it is a different source/center for the same coefficient regime.  Thus,
train-vs-heldout gaps measure source generalization inside one coefficient
field, not broad coefficient generalization.

\section{Learned preconditioner families}

All learned models are intended to approximate or improve an action
\[
  z_\theta = B_\theta r
\]
used inside stationary iteration or FGMRES.  The operator may be nonlinear in
\(r\); when it is nonlinear or changes between applications, flexible GMRES is
the appropriate Krylov wrapper.

\subsection{U-Net residual model}

The residual U-Net first applies the chosen classical base preconditioner:
\[
  z_0=\Bh r.
\]
It feeds real and imaginary channels of \(z_0\), together with coefficient
features such as \(\sigma_t\) and \(\sigma_s\), into a gated U-Net correction:
\begin{equation}
  B_\theta r
  =
  z_0 + C_\theta(z_0,\sigma_t,\sigma_s).
\end{equation}
This model is residual-dependent through \(z_0\), but it does not explicitly
represent a second defect correction.

\paragraph{Coefficient-dependent multilevel gates.}
The U-Net correction path is deliberately linear in the residual features for
fixed coefficients; nonlinearity enters through coefficient-dependent gates.
The original implementation already used one coefficient gate at every encoder
and decoder level.  A coefficient encoder produced a feature pyramid
\(\{c_\ell\}\), and each residual feature map \(u_\ell\) was modulated as
\[
  u_\ell \leftarrow
  u_\ell\left(1+\gamma\tanh(P_\ell c_\ell)\right),
\]
where \(P_\ell\) is a level-specific \(1\times1\) convolution.  This gave
different learned gates at different U-Net levels, but the gate was mostly
local at the current resolution.

The updated \code{--scale-aware-gates} option makes the multilevel mechanism
more explicit.  At each level \(\ell\), the gate now uses the learned
coefficient feature \(c_\ell\), a raw physical coefficient feature
\(p_\ell\), and scale summaries.  The raw feature includes normalized
\(\sigma_t\), albedo, coefficient gradients, anisotropy \(g\), coarse-cell
optical thickness \(2^\ell h\sigma_t\), absorption ratio
\(\sigma_a/\sigma_t\), transport ratio \(\sigma_{\rm tr}/\sigma_t\), and
\(\sigma_{\rm tr}=\sigma_t-g\sigma_s\).  The local gate input is
\(\tilde c_\ell=[c_\ell,p_\ell]\).  Its summary is
\[
  s_\ell =
  \left[
    \operatorname{mean}(\tilde c_\ell),
    \operatorname{std}(\tilde c_\ell),
    \min(\tilde c_\ell),
    \max(\tilde c_\ell),
    \ell/(L-1),
    \log_2(2^\ell),
    2^{-\ell},
    \operatorname{stage}
  \right],
\]
\[
  u_\ell \leftarrow
  u_\ell\left(
    1+
    \gamma_\ell
    \tanh\{P_\ell \tilde c_\ell + M_\ell s_\ell\}
  \right).
\]
Here \(\operatorname{stage}=-1\) for encoder/restriction-side gates and
\(\operatorname{stage}=1\) for decoder/prolongation-side gates.  \(M_\ell\) is
a small MLP and \(\gamma_\ell\) is a bounded learnable gate strength.  The
optional \code{--scale-gate-channel-mix} further adds a bounded
coefficient-conditioned channel-mixing correction
\[
  u_\ell \leftarrow u_\ell
  +
  \alpha_\ell W_\ell(s_\ell)u_\ell ,
\]
initialized at zero.  For fixed coefficients this still defines a linear
operator in the residual, but it allows the gate to rotate or couple residual
feature channels rather than merely reweight them.  Thus fine levels can learn
coefficient dependence appropriate for oscillatory/local error components,
while coarse levels can learn dependence on coarser optical-thickness and
albedo summaries.  This is closer to the multilevel-smoother interpretation:
different U-Net scales are allowed to accelerate different error subspaces with
different coefficient conditioning.

\subsection{U-Net defect2 model}

The \code{UNet-defect2} model computes one explicit defect correction before
learning:
\begin{align}
  z_0 &= \Bh r,\\
  d &= r-\Ah z_0,\\
  z_1 &= \Bh d,
\end{align}
and returns
\begin{equation}
  B_\theta r
  =
  z_0+\omega z_1+C_\theta(z_0,z_1,\sigma_t,\sigma_s).
  \label{eq:unet-defect2}
\end{equation}
This is the strongest learned model observed so far when \(\Bh\) is
full-angular TSA.  Mechanistically, the model starts from two physically
meaningful residual-correction directions and lets the U-Net learn a nonlinear
defect update on top.

\subsection{Routed learned base inside U-Net defect2}

In the original \code{UNet-defect2} experiments, \(B_h\) is fixed by the
command-line \code{base_kind}.  For example, with
\code{--base-kinds transport_tsa},
\[
  B_h r = B_{\rm TSA}r,
\]
and the model computes \(z_0=B_{\rm TSA}r\) and
\(z_1=B_{\rm TSA}(r-A_hz_0)\).  With \code{--base-kinds transport_dsa},
the same formula uses \(B_{\rm DSA}\).  Thus the earlier model learned only the
defect correction \(C_\theta\), not the base preconditioner \(B_h\).

The updated ablation code includes a \code{LearnedBaseRouter} option.  In this
variant,
\[
  B_{\theta}^{\rm route}r
  =
  \sum_{j=1}^{J}\pi_j(\mu,r)\,
  \left[
    \alpha_j(\mu,r)G_\eta r
    +
    \beta_j(\mu,r)S_j(r-A_hG_\eta r)
  \right],
  \label{eq:routed-base}
\]
where \(\mu\) denotes coefficient features derived from
\(\sigma_t,\sigma_s,g,h\), and \(S_j\) is one of the synthetic corrections:
zero for the pure transport candidate, scalar diffusion for DSA, full
reference transport-scattering for TSA, or closed moment TSA.  The route
weights \(\pi_j\) are produced by a softmax router, while
\(\alpha_j,\beta_j\) are bounded learned relaxation factors.  The router is
initialized close to the selected classical \code{base_kind}, so the initial
map remains near a known stable preconditioner.

This design addresses two limitations of the earlier \code{UNet-defect2}
model.  First, it can choose between transport-, DSA-, TSA-, and moment-TSA-like
corrections according to the problem coefficients.  Second, it can tune the
relative transport and synthetic strengths in a coefficient-conditioned way.
The router can optionally use residual summaries, making \(B_\theta^{\rm route}\)
nonlinear in the residual.  In that case FGMRES, not ordinary GMRES, is the
appropriate Krylov wrapper.

One subtle point is the loss.  If \(B_h\) itself is learned, then a metric loss
of the form \(\norm{B_h r^+}/\norm{B_h r}\) can be gamed by changing the metric.
The implementation therefore keeps the diagnostic metric tied to the fixed
classical base while the routed \(B_\theta^{\rm route}\) is used in the actual
correction.  Pure physical loss remains the cleanest first objective for routed
models.

\subsection{Moment U-Net}

\code{MomentUNet} applies the same defect2 idea but restricts the learned
correction to \(\operatorname{range}(\QK)\).  It is safer and cheaper than
full-angular learning, but it can miss high angular modes in anisotropic and
mixed regimes.

\subsection{Coefficient-weighted TSA}

\code{CoeffWeightedTSA-MK} predicts bounded spatial weights for retained
angular Fourier magnitudes and rescales the fixed moment TSA correction:
\begin{equation}
  \widehat z_m(x)\leftarrow w_{|m|,\theta}(x)\widehat z_m(x),
  \qquad |m|\le K.
\end{equation}
This architecture has about \(2.27\times10^4\) parameters in the width-24
\(K=4\) runs, but the main limitation is structural rather than parameter
count.  It cannot rotate modes, couple modes, learn restriction/prolongation,
or choose an angular subspace based on the residual.  It only rescales a fixed
closed-form moment correction.

\subsection{Coefficient-mixed TSA}

\code{CoeffMixedTSA-MK} predicts a local matrix \(M_\theta(x)\) acting on the
retained moment coefficients:
\begin{equation}
  \widehat{\bm z}(x)\leftarrow M_\theta(x)\widehat{\bm z}(x).
\end{equation}
This gives mode coupling inside the fixed \(\QK\) space.  It is more expressive
than diagonal coefficient weighting, but it still cannot leave the prescribed
angular Fourier subspace.

\subsection{Learned restriction/prolongation TSA}

\code{LearnedRP-TSA-MK} predicts coefficient-conditioned restriction and
prolongation maps around the closed moment inverse:
\begin{equation}
  \QK^*s \mapsto R_\theta(x)\QK^*s,\qquad
  \widehat{\bm z}(x) \mapsto P_\theta(x)\widehat{\bm z}(x).
\end{equation}
This is the most promising structured moment candidate because it modifies both
the low-order right-hand side and the lifted correction.  Current
implementations can preserve the scalar moment and are initialized near the
identity.  However, they still operate inside the fixed \(\QK\) candidate
space, so they are not yet true learned angular subspaces.

The identity anchor used in bottleneck experiments is
\begin{equation}
  L_{\rm id}
  =
  \frac12\norm{R_\theta-I}_F^2
  +
  \frac12\norm{P_\theta-I}_F^2.
\end{equation}
It improves FGMRES safety but often suppresses learning, selecting models close
to the initial stable classical preconditioner.

\subsection{Residual-dependent subspace TSA}

\code{ResidualSubspaceTSA-MK} predicts a global matrix from summaries of the
current residual moments:
\begin{equation}
  S_\theta(r)
  =
  I+\Delta_\theta\left(\operatorname{summary}(\QK^*s)\right).
\end{equation}
This is genuinely nonlinear in the residual and can in principle adapt angular
directions to the current error.  The current experiments show useful one-step
movement but poor FGMRES stability on the mixed case, so norm control,
near-identity continuation, and Krylov-aware training are required.

\section{Diagnostics and training objectives}

For any preconditioner \(B_\theta\), one application gives
\begin{equation}
  z_\theta=B_\theta r,\qquad r^+=r-\Ah z_\theta.
\end{equation}
The three main one-step ratios are
\begin{align}
  \rho_{\rm phys}(r)
  &=
  \frac{\norm{r^+}_2}{\norm r_2},
  \label{eq:rho-phys}\\
  \rho_{\rm born}(r)
  &=
  \frac{\norm{\Geta r^+}_2}{\norm{\Geta r}_2},\\
  \rho_{\rm base}(r)
  &=
  \frac{\norm{\Bh r^+}_2}{\norm{\Bh r}_2}.
\end{align}
The field \code{one_step_true} in the CSV files is \(\rho_{\rm phys}\).  The
heldout physical loss is therefore
\begin{equation}
  L_{\rm phys}^{\rm heldout}
  =
  \rho_{\rm phys}(r_{\rm heldout})^2
  =
  \code{one_step_true}^2.
\end{equation}

This is not the same object as \code{train_loss_final}.  The saved training
loss is a sampled objective over residual batches and may include metric-space
ratios, physical ratios, correction-size regularization, rollout terms,
stability barriers, teacher imitation, and identity penalties.  A large drop in
\code{train_loss_final} means the sampled objective decreased; it does not
guarantee that the train-source \(\rho_{\rm phys}\), heldout
\(\rho_{\rm phys}\), or FGMRES residual improved.

The general loss in \path{rte/npbs_unet_ablation.py} can be written as
\begin{equation}
  L
  =
  \lambda_m\mathbb{E}\rho_{\rm metric}^2
  +
  \lambda_p\mathbb{E}\rho_{\rm phys}^2
  +
  \lambda_c\mathbb{E}\rho_{\rm corr}^2
  +
  L_{\rm rollout}
  +
  L_{\rm stab}
  +
  L_{\rm teacher}.
  \label{eq:general-loss}
\end{equation}
Here \(\rho_{\rm metric}\) is usually a Born- or base-preconditioned residual
ratio, not the raw physical residual.  The correction-size term has the form
\begin{equation}
  \rho_{\rm corr}
  =
  \frac{\norm{B_\theta r}_2}{\norm{M r}_2},
\end{equation}
where \(M\) is the metric map chosen in the harness.  Rollout terms repeatedly
apply the learned correction:
\begin{equation}
  r_{j+1}=r_j-\Ah B_\theta r_j,\qquad
  L_{\rm rollout}
  =
  \lambda_{\rm roll}
  \sum_{j=1}^{J}\gamma^{j-1}\rho(r_j)^2 .
\end{equation}
The stability barrier penalizes ratios above a target:
\begin{equation}
  L_{\rm stab}
  =
  \lambda_{\rm stab}\frac1J
  \sum_j
  \max(\rho(r_j)-\rho_\star,0)^2 .
\end{equation}
Teacher imitation is
\begin{equation}
  L_{\rm teacher}
  =
  \lambda_T\mathbb{E}
  \frac{\norm{B_\theta r-Tr}_2^2}{\norm{Tr}_2^2},
\end{equation}
where \(T\) is a stronger classical correction such as depth-4 moment TSA or
full TSA.

In \path{rte/npbs_bottleneck_harness.py}, \code{loss_mode=raw} means pure
physical one-step training.  \code{loss_mode=raw_rollout_teacher} means raw
physical one-step plus physical rollout and teacher imitation.  The
\code{metric_map} option can be \code{base}, \code{forward}, or
\code{damped_moment}.  For structured nonlinear models, the safest reported
selection criterion so far is not final sampled training loss, but a composite
view including heldout physical loss and heldout FGMRES.

\section{Why loss can decrease but one-step or FGMRES may not improve}

The main observed failure mode is objective mismatch.  A learned model can
reduce the recorded training loss while failing to become a better solver
preconditioner for several reasons.

\begin{enumerate}
  \item The saved training loss may be a metric-space ratio rather than the
  raw physical residual \(\norm{r-\Ah B_\theta r}/\norm r\).
  \item The residual batches include random and rollout residuals, not only the
  train source used in the table's \code{split=train} one-step diagnostic.
  \item One-step contraction is necessary but not sufficient for Krylov
  quality.  The learned \(B_\theta\) can reduce \(\norm{r^+}\) for sampled
  residuals while producing a preconditioned operator with worse spectrum or
  worse non-normal transient growth.
  \item Nonlinear residual-dependent maps should be used with FGMRES.  Even
  with FGMRES, aggressive residual-dependent corrections can create unstable
  local linearizations.
  \item Moment-space models may reduce the component visible inside \(\QK\)
  while leaving important angular error outside \(\QK\).  This is especially
  severe in anisotropic and mixed cases.
\end{enumerate}

This explains why the main \code{UNet-defect2} models can train strongly in
sampled loss, while structured moment models can show one-step improvement but
still degrade FGMRES.

\section{Current numerical evidence}

All numbers below are from committed CSV artifacts in \path{rte}.  The
``heldout loss'' column is computed as \(\code{one_step_true}^2\).

\subsection{Full-angular TSA plus U-Net defect2}

The strongest positive learned result is the physical-loss full-angular
\code{UNet-defect2} run
\path{rte/npbs_unet_ablation_tsa_full_physical_steps900.csv}.  It uses
\code{transport_tsa} as the base and pure physical one-step loss.

\begin{center}
\small
\begin{tabular}{llrrrr}
\toprule
Case & Method & Train one-step & Heldout one-step & Heldout loss & Heldout FGMRES \\
\midrule
anisotropic & \code{tsa2} & -- & 0.08236 & 0.00678 & \(1.23\cdot10^{-5}\) \\
anisotropic & \code{tsa4} & -- & 0.03645 & 0.00133 & \(1.30\cdot10^{-5}\) \\
anisotropic & \code{UNet-defect2-w1} & 0.08609 & 0.06990 & 0.00489 & \(1.27\cdot10^{-5}\) \\
\midrule
mixed & \code{tsa2} & -- & 0.17898 & 0.03203 & \(7.60\cdot10^{-6}\) \\
mixed & \code{tsa4} & -- & 0.11446 & 0.01310 & \(7.94\cdot10^{-6}\) \\
mixed & \code{UNet-defect2-w1} & 0.16587 & 0.12453 & 0.01551 & \(7.54\cdot10^{-6}\) \\
\bottomrule
\end{tabular}
\end{center}

For the mixed case, the U-Net almost matches \code{tsa4} in heldout one-step
loss and slightly improves the final FGMRES residual in this run.  This is the
best evidence that learning can help when the base representation contains
full angular information.

\subsection{Structured moment-space models}

The moment-structure physical-loss run
\path{rte/npbs_unet_ablation_moment_structure_physical_size64_steps1800.csv}
uses \(64^2\) grids, 24 angles, \(K=4\) moments, and relaxation \(0.4\) for the
moment-TSA base.  The mixed case is the most informative.

\begin{center}
\small
\begin{tabular}{lrrrrr}
\toprule
Method & Train one-step & Heldout one-step & Heldout loss & FGMRES & Sampled loss \\
\midrule
\code{tsa_moment1-w0.4} & -- & 0.64758 & 0.41936 & \(1.97\cdot10^{-5}\) & -- \\
\code{tsa_moment4-w0.4} & -- & 0.24397 & 0.05952 & \(7.32\cdot10^{-6}\) & -- \\
\code{CoeffWeightedTSA-M4} & 0.64788 & 0.60235 & 0.36283 & \(1.93\cdot10^{-5}\) & \(0.6035\to0.5894\) \\
\code{CoeffMixedTSA-M4} & 0.57202 & 0.59097 & 0.34925 & \(4.29\cdot10^{-5}\) & \(0.6014\to0.5786\) \\
\code{LearnedRP-TSA-M4} & 0.51855 & 0.50000 & 0.25000 & \(3.21\cdot10^{-4}\) & \(0.5975\to0.5800\) \\
\code{ResidualSubspaceTSA-M4} & 0.58649 & 0.58787 & 0.34559 & \(1.23\cdot10^{-4}\) & \(0.6042\to0.5886\) \\
\bottomrule
\end{tabular}
\end{center}

The interpretation is sharp.  Diagonal coefficient weighting barely moves.
Mode mixing and learned restriction/prolongation reduce the one-step residual
on both train and heldout sources, so mode-direction learning is not a dead
idea.  However, FGMRES becomes much worse for the more aggressive learned maps,
especially \code{LearnedRP-TSA-M4}.  This is not ordinary source overfitting:
train and heldout one-step numbers are similar.  The bottleneck is stability
and solver alignment.

For the anisotropic case in the same file, FGMRES is less sensitive, but the
pattern remains: learned R/P and residual subspace models improve one-step
loss over \code{tsa_moment1-w0.4}, yet do not approach the deeper
\code{tsa_moment4-w0.4} one-step correction.

\begin{center}
\small
\begin{tabular}{lrrrr}
\toprule
Method & Train one-step & Heldout one-step & Heldout loss & FGMRES \\
\midrule
\code{tsa_moment1-w0.4} & -- & 0.65593 & 0.43025 & \(1.27\cdot10^{-5}\) \\
\code{tsa_moment4-w0.4} & -- & 0.26037 & 0.06779 & \(1.25\cdot10^{-5}\) \\
\code{CoeffWeightedTSA-M4} & 0.62549 & 0.62493 & 0.39054 & \(1.29\cdot10^{-5}\) \\
\code{CoeffMixedTSA-M4} & 0.60338 & 0.60656 & 0.36791 & \(1.24\cdot10^{-5}\) \\
\code{LearnedRP-TSA-M4} & 0.56925 & 0.57530 & 0.33097 & \(1.30\cdot10^{-5}\) \\
\code{ResidualSubspaceTSA-M4} & 0.58391 & 0.59881 & 0.35858 & \(1.25\cdot10^{-5}\) \\
\bottomrule
\end{tabular}
\end{center}

\subsection{Checkpointing and identity anchoring}

The bottleneck harness tested learned R/P with teacher and rollout losses on
the mixed size-64 case.  An uncheckpointed run improves one-step residual but
hurts FGMRES.  Checkpointing by heldout FGMRES and adding an identity anchor
restores solver safety, but mostly selects a near-initial operator.

\begin{center}
\small
\begin{tabular}{lrrrrr}
\toprule
Run & Selected step & Heldout one-step & Heldout loss & FGMRES & Sampled loss \\
\midrule
uncheckpointed teacher 0.5
& final & 0.50156 & 0.25156 & \(1.91\cdot10^{-4}\) & \(1.1148\to1.0781\) \\
cosine checkpoint
& 1 & 0.64730 & 0.41900 & \(1.96\cdot10^{-5}\) & \(1.1243\to1.0778\) \\
onecycle checkpoint
& 1 & 0.64756 & 0.41933 & \(1.96\cdot10^{-5}\) & \(1.1196\to1.0740\) \\
anchor \(5\cdot10^{-3}\), checkpoint
& 100 & 0.62394 & 0.38931 & \(1.91\cdot10^{-5}\) & \(1.1367\to1.1045\) \\
anchor \(10^{-2}\), checkpoint
& 100 & 0.62497 & 0.39059 & \(1.94\cdot10^{-5}\) & \(1.1260\to1.0984\) \\
\bottomrule
\end{tabular}
\end{center}

This table is important for deciding the next design.  The problem is not that
the model cannot change the one-step residual.  The problem is that the learned
change is not constrained to remain a good Krylov preconditioner.

\section{Current bottleneck diagnosis}

The bottlenecks, in order of importance, are as follows.

\begin{enumerate}
  \item \textbf{Fixed angular subspace.}
  \code{CoeffWeightedTSA-M4} is small, but the decisive issue is that it cannot
  learn better modes.  It only rescales fixed moment corrections.  Even
  \code{CoeffMixedTSA} and current \code{LearnedRP-TSA} stay inside fixed
  \(\QK\).

  \item \textbf{Lack of stable residual-dependent angular directions.}
  Hard mixed and anisotropic cases require mode directions that depend on
  coefficient regime and residual content.  \code{ResidualSubspaceTSA} begins
  to address this, but its unconstrained residual dependence can destabilize
  FGMRES.

  \item \textbf{One-step loss is not solver aligned.}
  Physical loss is clearer than metric loss, but even physical one-step
  residual does not control the spectrum or non-normality of \(I-\Ah B_\theta\).

  \item \textbf{Checkpoint criteria are too late or too conservative.}
  FGMRES checkpointing catches instability but tends to select near-initial
  models.  A better training objective should make safe movement possible
  rather than merely rejecting unsafe movement.

  \item \textbf{No separately saved heldout sampled objective.}
  Current CSVs report train sampled loss and train/heldout one-step source
  diagnostics.  They do not always report the same sampled objective evaluated
  on a heldout residual bank.  This makes it harder to separate optimization,
  source generalization, and solver-alignment failures.
\end{enumerate}

\section{Recommended next architectures}

\subsection{Candidate-space learned restriction/prolongation}

The highest-priority design is a true candidate-space learned R/P model.  Use a
larger angular candidate basis \(Q_{K_{\rm cand}}\) with
\(K_{\rm cand}>K_{\rm solve}\), for example \(K_{\rm solve}=4\) and
\(K_{\rm cand}=6\) or \(8\).  Learn
\begin{equation}
  R_\theta(x,r): \operatorname{range}(Q_{K_{\rm cand}})\to
  \operatorname{range}(Q_{K_{\rm solve}}),
  \qquad
  P_\theta(x,r): \operatorname{range}(Q_{K_{\rm solve}})\to
  \operatorname{range}(Q_{K_{\rm cand}}),
\end{equation}
then lift back to the full angular grid through \(Q_{K_{\rm cand}}\).  This
finally lets the model rotate and combine angular directions outside the fixed
solve space while retaining a small closed moment inverse.

The model should preserve the scalar mode exactly or nearly exactly:
\[
  R_\theta e_0=e_0,\qquad P_\theta e_0=e_0,
\]
because scalar flux is the dominant stable low-order quantity in diffusive
transport.  It should also enforce conjugate symmetry of angular Fourier
coefficients so the lifted angular field remains physically consistent.

\subsection{Bounded low-rank updates}

Instead of predicting dense unrestricted \(R_\theta\) and \(P_\theta\), use
near-identity bounded low-rank updates:
\begin{align}
  R_\theta &= I + \alpha\, U_\theta V_\theta^*,\\
  P_\theta &= I + \alpha\, \widetilde U_\theta \widetilde V_\theta^*,
\end{align}
with \(\alpha\) learned or scheduled but bounded.  Spectral or Frobenius norm
penalties should be applied to the update, not only to the final residual.
This gives the model a controlled way to rotate modes without destroying the
classical preconditioner's stability.

\subsection{Residual-gated angular subspace}

A residual-dependent gate should be retained, but made conservative:
\begin{equation}
  B_\theta r
  =
  B_{\rm safe}r
  +
  g_\theta(r,\sigma_t,\sigma_s)
  \, \Delta B_\theta r,
  \qquad
  0\le g_\theta\le g_{\max}.
\end{equation}
Here \(B_{\rm safe}\) can be damped \code{tsa_moment1-w0.4} or
\code{tsa_moment4-w0.4}; \(\Delta B_\theta\) is the learned candidate-space R/P
correction.  The gate should use residual moment summaries, not only
coefficient features, because the useful angular subspace depends on the
current error.

\subsection{Spectrum-aware and Krylov-aware losses}

The next training objective should combine raw physical one-step loss with
terms that discourage bad preconditioned spectra:
\begin{equation}
  L
  =
  L_{\rm phys}
  +
  \lambda_{\rm roll}L_{\rm rollout}
  +
  \lambda_TL_{\rm teacher}
  +
  \lambda_{\rm grow}L_{\rm growth}
  +
  \lambda_{\rm id}L_{\rm id}.
\end{equation}
Useful growth proxies include random probes of the error propagator
\[
  E_\theta v = v-\Ah B_\theta v,
\]
for example
\[
  L_{\rm growth}
  =
  \mathbb{E}_v\max\left(
  \frac{\norm{E_\theta v}}{\norm v}-\tau,0
  \right)^2
  +
  \mathbb{E}_v\max\left(
  \frac{\norm{E_\theta^2 v}}{\norm{E_\theta v}}-\tau_2,0
  \right)^2 .
\]
For small laptop tests, an even more direct objective is a short unrolled
FGMRES or Anderson-like loss on a few residuals.  It will be slower, but it
targets the actual failure observed in the tables.

\section{Recommended experimental protocol}

Every new learned preconditioner should report at least the following fields:
\begin{itemize}
  \item train sampled loss initial/final;
  \item heldout sampled loss initial/final, using a heldout residual bank;
  \item train-source \(\rho_{\rm phys}\) and \(\rho_{\rm phys}^2\);
  \item heldout-source \(\rho_{\rm phys}\) and \(\rho_{\rm phys}^2\);
  \item \(\rho_{\rm born}\) and \(\rho_{\rm base}\);
  \item stationary residual history when meaningful;
  \item FGMRES residual history, final residual, and time;
  \item base-call count and wall-clock time;
  \item parameter count, architecture dimensions, moment cutoffs, and whether
  the model is nonlinear in the residual;
  \item checkpoint metric and selected step.
\end{itemize}

The minimum hard-case suite should include \code{anisotropic} and
\code{mixed} at size \(64\), 24 angles for moment-space models, and both
train/heldout sources.  The \code{mixed} case should be treated as the primary
stability test.

\section{Reproducible command templates}

The following commands are templates matching the current parser options.  They
may need device changes on a laptop.

\paragraph{Full-angular U-Net defect2 with physical loss.}
\begin{verbatim}
python -m rte.npbs_unet_ablation \
  --cases anisotropic mixed \
  --sizes 32 64 \
  --n-angles 12 \
  --base-kinds transport_tsa \
  --forms defect2 \
  --scale-aware-gates \
  --scale-gate-channel-mix \
  --defect-relaxations 1.0 \
  --baseline-depths 1 2 3 4 \
  --baseline-relaxations 1.0 0.5 \
  --steps 900 \
  --lr 0.002 \
  --scheduler step \
  --scheduler-step-size 300 \
  --scheduler-gamma 0.5 \
  --metric-weight 0.0 \
  --physical-weight 1.0 \
  --correction-weight 0.0 \
  --rollout-weight 0.0 \
  --fgmres-mode complex \
  --output-csv rte/npbs_unet_ablation_tsa_full_physical_steps900.csv \
  --output-md rte/NPBS_UNET_ABLATION_REPORT.md
\end{verbatim}

\paragraph{Routed learned-\(B_h\) U-Net defect2.}
\begin{verbatim}
python -m rte.npbs_unet_ablation \
  --cases anisotropic mixed \
  --sizes 64 \
  --n-angles 24 \
  --base-kinds transport_tsa \
  --forms defect2 \
  --scale-aware-gates \
  --scale-gate-channel-mix \
  --learned-base-router \
  --learned-base-candidates transport dsa tsa tsa_moment \
  --learned-base-use-residual \
  --learned-base-width 32 \
  --learned-base-blocks 2 \
  --learned-base-temperature 1.0 \
  --learned-base-route-scale 4.0 \
  --learned-base-relaxation-scale 0.5 \
  --moment-harmonics 4 \
  --steps 900 \
  --lr 0.002 \
  --scheduler cosine \
  --metric-weight 0.0 \
  --physical-weight 1.0 \
  --correction-weight 0.0 \
  --rollout-weight 0.0 \
  --fgmres-mode real-stacked \
  --output-csv rte/npbs_unet_ablation_routed_base_physical_steps900.csv \
  --output-md rte/NPBS_UNET_ABLATION_ROUTED_BASE_REPORT.md
\end{verbatim}

\paragraph{Moment-structure physical-loss comparison.}
\begin{verbatim}
python -m rte.npbs_unet_ablation \
  --cases anisotropic mixed \
  --sizes 64 \
  --n-angles 24 \
  --base-kinds transport_tsa_moment \
  --forms coeff_weighted_tsa coeff_mixed_tsa learned_rp_tsa residual_subspace_tsa \
  --moment-harmonics 4 \
  --moment-transport-relaxation 0.4 \
  --moment-synthetic-relaxation 0.4 \
  --baseline-depths 1 4 \
  --baseline-relaxations 0.4 \
  --steps 1800 \
  --lr 0.002 \
  --scheduler cosine \
  --width 24 \
  --levels 4 \
  --metric-weight 0.0 \
  --physical-weight 1.0 \
  --correction-weight 0.0 \
  --rollout-weight 0.0 \
  --output-csv rte/npbs_unet_ablation_moment_structure_physical_size64_steps1800.csv
\end{verbatim}

\paragraph{Learned R/P bottleneck harness with teacher, rollout, identity anchor,
and FGMRES checkpointing.}
\begin{verbatim}
python -m rte.npbs_bottleneck_harness \
  --case mixed \
  --size 64 \
  --n-angles 24 \
  --base-kind transport_tsa_moment \
  --model-kind learned_rp_moment_tsa \
  --moment-harmonics 4 \
  --loss-mode raw_rollout_teacher \
  --metric-map base \
  --teacher-depth 4 \
  --teacher-kind transport_tsa_moment \
  --teacher-harmonics 4 \
  --teacher-relaxation 0.4 \
  --teacher-weight 0.5 \
  --identity-weight 0.005 \
  --moment-transport-relaxation 0.4 \
  --moment-synthetic-relaxation 0.4 \
  --rollout-weight 0.2 \
  --rollout-steps 3 \
  --steps 1200 \
  --lr 0.0005 \
  --scheduler onecycle \
  --eval-every 100 \
  --select-checkpoint heldout_fgmres \
  --fgmres-mode real-stacked \
  --output-csv rte/npbs_bottleneck_harness_candidate.csv \
  --history-csv rte/npbs_bottleneck_harness_candidate_history.csv
\end{verbatim}

\section{Concrete next implementation tasks}

\begin{enumerate}
  \item Add heldout sampled-objective evaluation to
  \path{rte/npbs_bottleneck_harness.py} and \path{rte/npbs_unet_ablation.py}.
  This should evaluate the same loss components used for training on a heldout
  residual bank.

  \item Implement \code{CandidateRP-TSA}: \(K_{\rm solve}=4\),
  \(K_{\rm cand}=6\) and \(8\), scalar-mode preservation, conjugate symmetry,
  and bounded low-rank updates.

  \item Add residual-gated candidate-space updates:
  \(B_\theta=B_{\rm safe}+g_\theta\Delta B_\theta\), with a small maximum gate
  at the start and a continuation schedule.

  \item Add growth-probe diagnostics for \(E_\theta=I-\Ah B_\theta\):
  random one-step and two-step amplification, plus optional short unrolled
  FGMRES loss for very small batches.

  \item Use composite checkpointing.  A reasonable first rule is:
  accept checkpoints only if heldout FGMRES does not exceed the stable baseline
  by more than a tolerance, then choose the smallest heldout physical loss.

  \item Re-run the same \code{anisotropic} and \code{mixed} size-64 suite with
  24 angles, reporting train and heldout physical losses, heldout sampled
  losses, and FGMRES histories.
\end{enumerate}

\section{Main conclusions}

\begin{enumerate}
  \item The RTE extension of NPBS is mathematically natural in the periodic
  pseudo-spectral testbed: the shifted/reference transport inverse \(\Geta\)
  gives the same Born identity structure as in Helmholtz.

  \item Classical TSA is the correct base family for hard RTE cases.  Raw
  transport/Born is too weak, and DSA only targets scalar diffusive errors.
  Moment TSA gives the right bridge between cheap low-order structure and
  learnable angular subspaces.

  \item Full-angular \code{UNet-defect2} trained with physical loss is the best
  current learned model.  It nearly matches deeper full TSA on the mixed case,
  which shows that learned defect corrections can be useful when the base
  representation contains enough angular information.

  \item \code{CoeffWeightedTSA-M4} is not competitive because it cannot learn
  modes.  Its parameter count is small, but the more important limitation is
  expressivity: it only rescales fixed moment corrections.

  \item Learned mode coupling and learned R/P are promising but currently
  unstable.  They reduce one-step residuals on train and heldout sources, so
  the learned direction is not meaningless.  But they often worsen FGMRES,
  showing that the loss must control solver-relevant operator behavior.

  \item The next decisive experiment is a stabilized candidate-space learned
  R/P model with residual-dependent but bounded gates, scalar-mode
  preservation, spectrum/growth penalties, teacher guidance, and checkpointing
  by heldout Krylov behavior.
\end{enumerate}

\end{document}
